\section{\label{sec:defns}光前分布与准分布}

本工作聚焦味非单态非极化的准分布与光前 PDF。向极化准 PDF 的推广是直接的，且大部分讨论同样适用于赝分布。我不考虑味单态情形，它会引入与胶子分布的额外混合；该混合在准分布情形下单圈微扰论中的研究见文献~\cite{Wang:2017qyg}。

\subsection{光前 PDF}

我用 $f(\xi,\mu)$ 表示重正化的光前 PDF，其中重正化标度为 $\mu$，并定义 $\xi=k^+/P^+$。这里光前坐标为 $(x^+,x^-,\boldsymbol{x}_\mathrm{T})$，且 $x^\pm = (t\pm z)/\sqrt{2}$。一般地，重正化光前 PDF 可用裸光前 PDF $f^{(0)}(\xi)$ 写成
\begin{equation}
f(\xi,\mu) = \int_\xi^1 \frac{\mathrm{d}\zeta}{\zeta}{\cal
Z}\left(\frac{\xi}{\zeta},\mu\right)f^{(0)}(\zeta).
\end{equation}
裸 PDF 定义为 \cite{Collins:2011zzd}
\begin{align}
 {} &f^{(0)}(\xi) = \int_{-\infty}^\infty
\frac{\mathrm{d}\omega^-}{4\pi}
e^{-i\xi P^+\omega^-} \nonumber\\
{} & \quad \times \left\langle P \left|
T\,\overline{\psi}(0,\omega^-,\mathbf{0}_\mathrm{T})
W(\omega^-,0)\gamma^+\frac{\lambda^a}{2}\psi(0) \right|
P\right\rangle_\mathrm{C},
\end{align}
其中 $T$ 是时序算符，$\psi$ 是夸克场，$\lambda^a$ 是 SU(2) 味 Pauli 矩阵，下标 C 表示真空期望值已被减去（换言之，只包含连通贡献）。算符 $W(\omega^-,0)$ 是 Wilson 线，
\begin{equation}
W(\omega^-,0) =
{\cal P}\exp\left[-ig_0\int_0^{\omega^-}\mathrm{d}y^-A^+_c(0,y^-,
\mathbf{0}_{\mathrm{T}})T_c\right],
\end{equation}
其中 ${\cal P}$ 是路径排序算符，$g_0$ 是 QCD 裸耦合，$A^\alpha = A^\alpha_c T_c$ 是具有生成元 $T_c$ 的 $SU(3)$ 规范势（对色指标 $c$ 的求和隐含在内）。靶态 $|P\rangle$ 是自旋平均的精确动量本征态，具有相对论性归一化
\begin{equation}
\langle P' | P \rangle =
(2\pi)^32P^+\delta\left(P^+-P^{\prime\,+}\right)\delta^{(2)}
\left(\mathbf{P}_\mathrm{T} - \mathbf{P}'_\mathrm{T}\right).
\end{equation}

\subsection{涂抹准分布}

准分布非微扰计算的中心挑战之一，是在格点调节子存在下与 Wilson 线长度相联系的幂次发散。已有多种移除该幂次发散的方案被提出 \cite{Chen:2016fxx,Ishikawa:2016znu,Chen:2017mzz,Alexandrou:2017huk,Orginos:2017kos}，包括用替代算符完全绕开这一问题 \cite{Aglietti:1998ur,Abada:2001if,Detmold:2005gg,Braun:2007wv,Bali:2017gfr}。在文献~\cite{Monahan:2016bvm} 中，我们通过梯度流 \cite{Narayanan:2006rf,Luscher:2011bx,Luscher:2013cpa} 同时涂抹费米子场与规范场，构造了从格点计算提取连续准分布的有限矩阵元。关于梯度流更完整的讨论，读者可参考近期综述 \cite{Luscher:2013vga,Ramos:2015dla}。

在有限流时下，梯度流保证由流场构造的矩阵元紫外有限。只要把流时以物理单位固定，就可以移除格点调节子，并提取作为流时函数的连续准分布。这一涂抹准分布随后可通过微扰匹配关系 \cite{Monahan:2016bvm} 与光前分布联系起来，类似于普通（未涂抹）准分布所满足的关系。这里，我在单圈微扰论水平研究定义涂抹准分布的矩阵元。该结果是把涂抹矩阵元与 $\ms$ 方案中的对应矩阵元联系起来所必需的。

连续涂抹矩阵元是有限的，但由于流时标度的存在，预期当流时趋于零时它包含幂次发散 \cite{Monahan:2016bvm}。该贡献必须非微扰地移除。与格点调节子不同，梯度流提供一个连续参数，这应当为研究幂次发散贡献提供更好的杠杆。在单圈微扰论中，该发散表现为与 Wilson 线长度线性的一项，第~\ref{sec:pertth} 节给出的结果证实了这一预期。

我把流时 $\tau$ 处的环标费米子场（ringed fermion fields）\cite{Makino:2014taa,Hieda:2016lly} 记为 $\overline{\chi}(x;\tau)$ 与 $\chi(x;\tau)$，把同一流时下由涂抹规范场 $B_\alpha(x;\tau)$ 构造的相应 Wilson 线记为 ${\bf\cal W}(x_1,x_2;\tau)$。我从连通矩阵元
\begin{align}\label{eq:hdef}
{} & \hm^{(s)}\left(\frac{n^2}{\tau},n\cdot P,
\sqrt{\tau}\Lambda_{\mathrm{QCD}},\sqrt{\tau}M_\mathrm{N}\right) = \nonumber\\
{} & \quad
\frac{1}{2P}\left\langle
P \left| \overline{\chi}(n;\tau)
{\bf\cal W}(n,0;\tau)\gamma_\alpha\frac{\lambda^a}{2}\chi(0;\tau)\right|
P\right\rangle_\mathrm{C},
\end{align}
出发，它是无量纲的，因而只依赖标度的无量纲组合。这里 $n$ 是通常取 $z$ 方向的 4-矢量，$n = (\mathbf{0},z,0)$；矩阵元对 4-矢量的依赖受欧氏 SO(4) 不变性约束 \cite{Radyushkin:2016hsy}。我指出流时的量纲是长度平方。环标费米子场不需要波函数重正化，并且只要流时 $\tau$ 非零且以物理单位固定，该涂抹矩阵元就是有限的，
因为由涂抹场构造的关联函数是有限的 \cite{Luscher:2011bx,Luscher:2013cpa}。发散将出现在流时趋于零的极限，届时矩阵元需要重正化。Lorentz 指标 $\alpha$ 可取 1 到 4 的任意值，但取 $\alpha = 4$ 可去除部分高扭度污染 \cite{Radyushkin:2016hsy}，并（至少对非极化味非单态分布）避免有限格距下的混合 \cite{Constantinou:2017sej}。在格点计算中，通常取核子态的空间动量为 $\mathbf{P} = (0,0,P_z)$，此时指标 $\alpha$ 限制为 $\alpha = \{3,4\}$。

涂抹准分布 \cite{Monahan:2016bvm} 在假设 $n = (\mathbf{0},z,0)$ 下定义为
\begin{align}\label{eq:qpdfdef}
q^{\,(s)}{} & \left(\xi,\sqrt{\tau}P_z,\sqrt{\tau}\Lambda_{\mathrm{QCD}},
\sqrt{\tau}M_\mathrm{N}\right)
=
\int_{-\infty}^\infty \frac{\mathrm{d}z}{2\pi} e^{i\xi z
P_z} P_z\nonumber \\
{} & \qquad \times h^{(s)}\left(\frac{n^2}{\tau},n\cdot P,\sqrt{\tau}\Lambda_{\mathrm{QCD}},
\sqrt{\tau}M_\mathrm{N}\right),
\end{align}
其中 $\xi$ 是一个无量纲参数，在欧氏空间中应将其视为（无量纲的）Fourier 共轭变量。

\subsection{涂抹准分布与 PDF 的联系}

涂抹准分布可通过 \cite{Monahan:2016bvm}
\begin{align}
{} & q^{\,(s)}\left(x,\sqrt{\tau} \Lambda_{\mathrm{QCD}}, \sqrt{\tau}P_z\right) =  \nonumber \\
{} & \; \int_{-1}^{1}
\frac{d\xi}{\xi}\, \widetilde{Z}\left(\frac{x}{\xi},\sqrt{\tau}\mu,
\sqrt{\tau}P_z\right)
f(\xi,\mu) +
{\cal O}(\sqrt{\tau}\Lambda_{\mathrm{QCD}})
\end{align}
直接与光前分布联系起来，条件是
\begin{equation}\label{eq:scales_tau}
\Lambda_{QCD},M_N\ll P_z \ll \tau^{-1/2}.
\end{equation}
准分布与光前 PDF 具有相同的红外结构 \cite{Ma:2014jla,Ma:2014jga,Carlson:2017gpk,Briceno:2017cpo}，因此匹配核可以在微扰论中确定。实践中，最简单的做法是先把涂抹准分布与零流时的准分布联系起来，再把这一（未涂抹的）准分布与光前 PDF 匹配 \cite{Xiong:2013bka,Ji:2015jwa}。

我在广义 Feynman 规范（取 $\alpha = \lambda = 1$，其中 $\alpha$ 是通常的规范固定参数，$\lambda$ 是为涂抹规范场固定规范而引入的参数 \cite{Luscher:2010iy,Luscher:2011bx}）下，分别在图~\ref{fig:sqpdf_tree} 与图~\ref{fig:sqpdf1} 中展示涂抹准分布的树图级与单圈贡献图形。这些图形对应单圈微扰论中的``夸克中的夸克''（quark-in-quark）分布。虽然准分布在轴规范下更容易计算，但该规范无法一致地推广到任意流时。
\begin{figure}
\centering
\includegraphics[width =.25\textwidth,keepaspectratio=true]{images/tree.pdf}
\caption{涂抹非单态准分布的树图级贡献。灰圈表示有限流时 $\tau$ 的费米子场，实线表示费米子传播子。}
\label{fig:sqpdf_tree}
\end{figure}
\begin{figure}
\centering
\includegraphics[width =0.4\textwidth]{images/qmatch1.pdf}\\
(a)\hspace*{0.12\textwidth}(b)\hspace*{0.12\textwidth}(c)\\[8pt]
\includegraphics[width =0.42\textwidth]{images/qmatch2.pdf}\\
(d)\hspace*{0.12\textwidth}(e)\hspace*{0.12\textwidth}(f)\\[8pt]
\includegraphics[width =0.48\textwidth]{images/qmatch_zpsi.pdf}\\
(g)\hspace*{0.1\textwidth}(h)\hspace*{0.1\textwidth}(i)\hspace*{0.1\textwidth}(j)
\caption{涂抹非单态准分布单圈贡献的图形。灰圈表示有限流时 $\tau$ 的费米子场；实线表示费米子传播子；双线是费米子流核；空心方块是任意流时 $\tau_1$ 处的流顶点。标记为 (a)、(d)、(f)、(g) 的图形与 $\ms$ 方案下未涂抹准分布的图形一一对应，后者中扩展算符里出现的费米子场换成了零流时的费米子场。}
\label{fig:sqpdf1}
\end{figure}
原则上，可以对这九个图求值并在单圈微扰论中直接确定匹配核。但实际上相应积分无法解析求值，更简单的做法是采用与文献~\cite{Ishikawa:2016znu,Constantinou:2017sej} 类似的程序，在固定的 Wilson 线长度 $z$ 处匹配。这一程序对应于在矩阵元本身层面（而非处理准分布）进行匹配，其更可取的主要原因在于矩阵元不依赖 $\xi$。在本工作的其余部分，我聚焦于矩阵元 $\hm^{(s)}$ 本身。
